\documentclass[journal]{IEEEtran}
\usepackage[T1]{fontenc}
\usepackage{graphicx}
\usepackage{booktabs}
\usepackage{amsmath}
\usepackage{array}
\usepackage{longtable}
\usepackage{tabularx}
\usepackage{enumitem}
\usepackage[hidelinks]{hyperref}
\usepackage{url}
\graphicspath{{../figures/}}
\newcommand{\figplace}[2]{\IfFileExists{figures/#1.png}{\includegraphics[width=\linewidth]{#1}}{\fbox{\parbox[c][0.22\textheight][c]{0.94\linewidth}{\centering #2}}}}
\newcolumntype{Y}{>{\raggedright\arraybackslash}X}
\begin{document}
\title{Multi-View and Multimodal Perception for Class-Wise Fruit Inventories: A Design-Space Review}
\author{Muhammad~Zainal~Muttaqin and Fatma~Indriani%
\thanks{Department of Computer Science, Universitas Lambung Mangkurat, Banjarbaru, Indonesia.}}
\maketitle
\begin{abstract}
Turning many observations of a plant into a single, class-wise fruit inventory
is not solved by per-image detection alone. A system must separate crowded
instances, attach a class label to each instance, and reject duplicates when the
same fruit reappears from another view. This design-space review maps verified
literature on RGB detection, depth estimation, multimodal fusion,
three-dimensional geometry, and cross-view association onto that task. Its domain
is agricultural fruit perception, with oil-palm fresh fruit bunches as the
principal case, but mechanisms are drawn freely from non-agricultural imaging,
including tracking, re-identification, multi-view detection, and structure from
motion, which enter as transferable evidence. Classification is treated as an
instance-dependent attribute, of which maturity grading is one case, because a
per-class count is undefined until the carrying instance has been identified.
The review synthesizes 182 full-text primary sources and extends the corpus with
a reproducible Scopus and OpenAlex search using seven query families, explicit
date bounds, and DOI-based deduplication. The search produced 21,066 resolved
records; 20,035 advanced after abstract screening; 250 were shortlisted; and 44
targeted full-text studies supplied verified evidence rows. The synthesis
distinguishes six identity mechanisms spanning intra-view detection, statistical
correction, appearance matching, temporal tracking, geometric association, and
learned multi-view association. A multi-scale real-time RGB detector is
identified as an appropriate baseline, depth as a conditional geometric cue
whose interpretation must be declared, and quality-aware fusion as a testable
alternative rather than an established agricultural result. No calibrated
depth-sensor benchmark for oil-palm bunches exists in the verified corpus. The
contribution is a testable design space for duplicate-aware, class-wise fruit
inventories and an explicit boundary between search candidates and verified
evidence.
\end{abstract}
\begin{IEEEkeywords}
multi-view perception, multimodal fusion, object detection, fruit counting,
instance identity, oil palm, design-space review
\end{IEEEkeywords}
\input{main5-body}
\bibliographystyle{IEEEtran}
\bibliography{references3}
\end{document}
